\section{Introduction}

Before launching a campaign, marketers want to know which version of a message will land. A fast-
growing practice replaces (or precedes) live A/B testing with \emph{synthetic audience
simulation}: a large language model (LLM) is conditioned on a set of audience ``personas'' and
asked to react to each candidate message, and the message its personas prefer is shipped. The
appeal is obvious---instant, cheap feedback---and it is encouraged by findings that
profile-conditioned LLMs can reproduce aspects of real human samples, an idea sometimes called
``silicon sampling''~\citep{argyleSiliconSampling}.

The appeal, however, outruns the evidence. The premise that a synthetic persona predicts how a
\emph{real} audience behaves is rarely tested against real outcomes, because doing so requires
ground truth that pairs concrete copy with measured audience response. We borrow the
\emph{sim-to-real} framing from robotics---how well does behaviour measured in simulation transfer
to the real world?---and instantiate it with the Upworthy Research
Archive~\citep{matiasUpworthy}: thousands of A/B tests (which we call \emph{packages}) in which
competing headline variants were shown to the same real traffic and their click-through rates
recorded. This lets us ask, for held-out copy never seen during persona design, whether the
simulator orders the variants the way the audience did.

We report two main findings. First, the dominant obstacle is not the model but the data: in most
A/B tests the real winner is not statistically distinguishable from the runner-up, so no predictor
can recover it---validity is only meaningful on the reliable subset. Second, and contrary to the
premise of persona simulation, \emph{persona conditioning hurts}: a no-persona zero-shot baseline
that simply asks the model how likely a typical reader is to click ranks the variants markedly
better (a medium effect) than the demographically grounded ten-persona panel (a small,
barely-significant effect), with non-overlapping confidence intervals. The base model holds an
accurate population-level prior on clickability; forcing it to role-play specific personas pulls
it away from that prior. We pre-register effect-size thresholds and report seed-stability and a
model comparison so neither effect is overstated.

\paragraph{Contributions.}
\begin{enumerate}[leftmargin=*]
  \item A reproducible \emph{sim-to-real} validity protocol for copy simulation, using a public,
        held-out A/B-tested benchmark~\citep{matiasUpworthy} as ground truth.
  \item A ground-truth-reliability filtering step (significant-winner A/B tests) that prevents
        validating a predictor against statistical noise~\citep{kohavi2009controlled}.
  \item A within-package ranking evaluation with a construct-matched click-intent elicitation,
        a no-persona zero-shot baseline that isolates the value of persona conditioning,
        pre-registered effect-size thresholds~\citep{cohen1988power}, bootstrap CIs, seed-stability,
        and a cost/quality model comparison.
  \item An artifact-first replication package: every reported number is regenerated from
        aggregated artifacts by public scripts.
\end{enumerate}

\paragraph{Research questions.}
\begin{description}[leftmargin=*]
  \item[RQ1] Do the personas' predicted ranking match the real click-through ranking on copy with
        a reliable winner, and do they beat a no-persona zero-shot baseline?
  \item[RQ2] Does predictive validity increase with the strength of the real effect?
  \item[RQ3] Is the result robust to the intent-vs-engagement construct gap, to seed sampling,
        and to model choice?
\end{description}
